\documentclass[11pt]{article}

\usepackage[margin=1in]{geometry}
\usepackage[T1]{fontenc}
\usepackage{lmodern}
\usepackage{xcolor}
\usepackage{amsmath}

\definecolor{navy}{HTML}{18324B}
\definecolor{teal}{HTML}{167D8D}
\definecolor{lightblue}{HTML}{EEF5F7}
\definecolor{warm}{HTML}{FFF6E8}
\definecolor{muted}{HTML}{566573}

\setlength{\parindent}{0pt}
\setlength{\parskip}{7pt}

\newcommand{\highlight}[1]{%
  \begin{center}
  \fcolorbox{teal}{lightblue}{\parbox{0.90\textwidth}{\centering\large\bfseries\color{navy}#1}}
  \end{center}}

\begin{document}

{\color{teal}\small DECISION INTELLIGENCE PROTOTYPE}\par\vspace{5pt}
{\color{navy}\fontsize{26}{30}\selectfont\bfseries Capacity allocation\par}
{\color{navy}\fontsize{26}{30}\selectfont\bfseries for an IAM operations queue\par}
\vspace{8pt}
{\color{muted}\large A plain-English look at choosing capacity when demand is uncertain.\par}
\vspace{14pt}

\highlight{How much hourly capacity should an IAM operations system keep available when user demand is uncertain and scheduled work is known in advance?}

\section*{Why this question matters}

Many operational systems have two kinds of work. User-driven requests arrive
unpredictably. Other work is scheduled in advance: synchronizations, access
reviews, reconciliations, exports, or policy updates.

Both types use the same people or infrastructure. If capacity is too low,
work waits and deadlines are missed. If capacity is too high, the system pays
for resources that may sit unused.

A forecast is useful, but it is not the decision. The decision also depends on
service expectations, backlog, deadlines, and cost.

\section*{The first model}

The prototype represents one shared hourly capacity pool:

\begin{itemize}
  \item user-driven demand is generated stochastically;
  \item user work is protected and served first;
  \item scheduled jobs use the capacity that remains; and
  \item unfinished work carries forward as backlog.
\end{itemize}

The current demonstration uses an hourly Poisson demand profile with lower
overnight activity and a daytime peak. Scheduled jobs have a workload, a
release time, and a deadline. Among eligible scheduled jobs, the simple
baseline processes the earliest deadline first.

\section*{The decision in one line}

For each candidate hourly capacity $c$, the model simulates many possible
demand paths and estimates the total cost:

\begin{equation*}
  \text{total cost} = \text{capacity cost}
  + \text{backlog cost}
  + \text{deadline penalties}.
\end{equation*}

The selected capacity is the candidate with the lowest estimated cost under
the stated assumptions.

\begin{center}
\renewcommand{\arraystretch}{1.35}
\begin{tabular}{p{0.27\textwidth} p{0.58\textwidth}}
\hline
\textbf{Model element} & \textbf{Meaning}\\
\hline
User demand & Uncertain work arriving during the hour\\
Capacity & Maximum work units that can be processed during the hour\\
Scheduled work & Known jobs competing for the remaining capacity\\
Backlog & Work that carries into a later hour\\
Cost & A way to express the trade-off between resources and waiting\\
\hline
\end{tabular}
\end{center}

\section*{What the experiment shows}

Under the current illustrative assumptions, the simulation selects an hourly
capacity of 17. This is not a recommendation for a real IAM deployment. It is
a reproducible example showing how a capacity decision can balance competing
consequences.

More capacity improves service and reduces backlog, but it also costs more.
Less capacity saves operating cost, but increases the chance that user work or
scheduled work waits. The best point depends on how those consequences are
valued.

\begin{center}
\fcolorbox{orange}{warm}{\parbox{0.88\textwidth}{\centering
\textbf{Important:} The data and cost weights in this first version are
synthetic and illustrative. The result demonstrates the mechanics of the
decision; it does not describe a production IAM environment.}}
\end{center}

\section*{What comes next}

The current version chooses the total capacity while keeping the scheduled-job
rule fixed. The next meaningful extension is to decide which scheduled jobs
should receive residual capacity, while continuing to protect user demand.

After that, the same decision can be compared with simpler rules, forecasting
models, and hybrid approaches using the same scenarios and evaluation metrics.

The purpose is not to add mathematics for its own sake. It is to test whether
additional structure produces a better operational decision.

\end{document}
